\section{Set-theoretic semantics}
\label{sec:semantics}

Our soundness framework is model-theoretic.  This is a deliberate
choice, and a departure from the existing rigorous treatment of the
CoqHammer translation: the soundness proof of~\cite{Czajka2016} for
the core translation is proof-theoretic --- it reconstructs a
type-theoretic proof term from a first-order proof in \emph{minimal
intuitionistic} first-order logic by induction on an $\eta$-long
normal derivation.  That technique does not extend to the questions
of this paper, for two reasons.  First, the ATPs actually invoked by
a hammer are \emph{classical}: a classical refutation is not a
minimal-logic derivation, so the reconstruction argument does not
even apply to the proofs the ATPs find.  Second, and decisively, the
optimisations we study --- monotonicity-based guard erasure above all
--- are justified by \emph{model constructions} (domain inflation,
sort merging), and, as \cite[\S6]{Czajka2016} itself observes, such
model-theoretic methods are incompatible with the proof-theoretic
soundness argument.  We therefore set up a set-theoretic semantics
for $\CICm$ once and for all, and measure every transformation
against it.  The price is that our soundness guarantee is semantic
(``the goal is true in the standard model'') rather than
proof-theoretic (``a proof term exists''); in the hammer setting this
is the appropriate guarantee, since the final arbiter is in any case
the proof-assistant kernel: the ATP output is used only to select
premises, and the goal is re-proved from them by certified search
\cite{CzajkaKaliszyk2018}.  What the semantic guarantee rules out is
the pipeline systematically wasting effort on, and reporting success
for, goals that are false; equivalently, it establishes that each
transformation adds no inconsistency to the generated problems.

Throughout, the metatheory is ZFC with two strongly inaccessible
cardinals; the axiom of choice is used and we flag the places where
it matters.

\begin{definition}[Standard model]
\label{def:model}
Fix Grothendieck universes $\Uone \in \Utwo$ (i.e.\ $\Uone =
V_{\kappa_1}$, $\Utwo = V_{\kappa_2}$ for inaccessibles $\kappa_1 <
\kappa_2$), and let $\Om = \{\emptyset, \holds\}$ where we write
$\holds$ for $\{\emptyset\}$; so $\Om = \{\emptyset, \{\emptyset\}\}
\in \Uone$.  Call a function $\rho$ from variables to sets a
\emph{$\Gamma$-assignment} if $\rho(x) \in \semr{A}{\rho}$ for every
$(x{:}A) \in \Gamma$ (well-defined by recursion on the context).  A
\emph{standard model} $\MM$ of an environment $\Env$ is a function
assigning to every judgement $\Gamma \vdash t : T$ and every
$\Gamma$-assignment $\rho$ a set $\semr{t}{\rho}$ (independent of the
derivation and of $T$; we omit $\rho$ for closed terms), such that:
\begin{enumerate}[label=(M\arabic*),leftmargin=3.2em]
\item \label{M:typing} \emph{(Typing soundness)}
  $\semr{t}{\rho} \in \semr{T}{\rho}$ whenever $\Gamma \vdash t : T$.
  Moreover $\semr{t}{\rho} = \emptyset$ whenever $t$ is a proof.
\item \label{M:sorts} \emph{(Sorts)}
  $\sem{\SProp} = \Om$, $\sem{\STypeI} = \Uone$,
  $\sem{\STypeII} = \Utwo$.
\item \label{M:var} \emph{(Variables, constants)}
  $\semr{x}{\rho} = \rho(x)$; for a definition $c \coloneqq t : T$,
  $\sem{c} = \sem{t}$.
\item \label{M:pi} \emph{(Products)}
  If $\Gamma, x{:}A \vdash B : \SProp$ then
  \[
    \semr{\Pi x{:}A.\,B}{\rho} \;=\;
    \setof{z \in \holds}{\forall a \in \semr{A}{\rho}.\;
        \semr{B}{\rho[x \mapsto a]} = \holds} .
  \]
  Otherwise (the product has sort $\STypeI$ or $\STypeII$),
  $\semr{\Pi x{:}A.\,B}{\rho}$ is the set-theoretic dependent
  product: the set of all functions (graphs) $f$ with domain
  $\semr{A}{\rho}$ such that $f(a) \in \semr{B}{\rho[x \mapsto a]}$
  for all $a \in \semr{A}{\rho}$.
\item \label{M:lam} \emph{(Abstraction, application)}
  If $\lambda x{:}A.\,t$ is not a proof, then
  $\semr{\lambda x{:}A.\,t}{\rho}$ is the graph
  $\setof{(a, \semr{t}{\rho[x \mapsto a]})}{a \in \semr{A}{\rho}}$;
  if $t\,u$ is not a proof, then
  $\semr{t\,u}{\rho} = \semr{t}{\rho}\bigl(\semr{u}{\rho}\bigr)$
  (function application of the graph).
\item \label{M:conv} \emph{(Conversion)}
  If $\Gamma \vdash t : T$, $\Gamma \vdash t' : T$ and $t \conv t'$,
  then $\semr{t}{\rho} = \semr{t'}{\rho}$.
\item \label{M:subst} \emph{(Substitutivity)}
  If $\Gamma, x{:}A, \Delta \vdash t : T$ and $\Gamma \vdash u : A$,
  then $\semr{t[u/x]}{\rho} = \semr{t}{\rho[x \mapsto \semr{u}{\rho}]}$
  for every assignment $\rho$ conforming to $\Gamma, \Delta[u/x]$.
\item \label{M:conn} \emph{(Connectives)}
  $\sem{\top} = \holds$, $\sem{\bot} = \emptyset$, and for
  $\Gamma$-propositions:
  \begin{align*}
  \semr{\varphi \wedge \psi}{\rho} &= \setof{z \in \holds}{\semr{\varphi}{\rho} = \holds \text{ and } \semr{\psi}{\rho} = \holds}, \\
  \semr{\varphi \vee \psi}{\rho} &= \setof{z \in \holds}{\semr{\varphi}{\rho} = \holds \text{ or } \semr{\psi}{\rho} = \holds}, \\
  \semr{\varphi \leftrightarrow \psi}{\rho} &= \setof{z \in \holds}{\semr{\varphi}{\rho} = \semr{\psi}{\rho}}, \\
  \semr{t \ceq_A u}{\rho} &= \setof{z \in \holds}{\semr{t}{\rho} = \semr{u}{\rho}}, \\
  \semr{\cexists{x}{A}\varphi}{\rho} &= \setof{z \in \holds}{\exists a \in \semr{A}{\rho}.\; \semr{\varphi}{\rho[x \mapsto a]} = \holds}.
  \end{align*}
\item \label{M:ind} \emph{(Inductive types)}
  For each inductive declaration as in \cref{def:env}(iii),
  $\sem{I}$ is a function $\Uone^p \to \Uone$, and for each $j$ the
  constructor denotation $\sem{c_j}$ is a (curried) function taking
  $\tup{a} \in \Uone^p$ and then argument tuples, such that for all
  $\tup{a} \in \Uone^p$:
  \begin{enumerate}[label=(\alph*)]
  \item \emph{(Closure and leastness)} $\sem{I}(\tup{a})$ is the
    least set $X \in \Uone$ such that for every $j$ and every tuple
    $\tup{d}$ with $d_l \in \sem{\sigma_{j,l}}_{[\tup{A} \mapsto
    \tup{a}]}$ for non-recursive positions and $d_l \in X$ for
    recursive positions, the value $\sem{c_j}(\tup{a})(\tup{d})$
    belongs to $X$.
  \item \emph{(Injectivity)} each $\sem{c_j}(\tup{a})$ is injective
    (jointly in all its value arguments).
  \item \emph{(Disjointness)} for $i \neq j$ the images of
    $\sem{c_i}(\tup{a})$ and $\sem{c_j}(\tup{a})$ on well-typed
    argument tuples are disjoint.
  \item \emph{(Exhaustiveness)} $\sem{I}(\tup{a})$ equals the union
    over $j$ of the images of $\sem{c_j}(\tup{a})$ on well-typed
    argument tuples.
  \end{enumerate}
\item \label{M:case} \emph{(Case)}
  If $\Gamma \vdash \case{I,\tup{\tau},Q}{t}{\tup{u}} : Q\,t$ is not
  a proof and $\semr{t}{\rho} = \sem{c_j}(\sem{\tup{\tau}}_\rho)(\tup{d})$
  (such $j$ and $\tup{d}$ exist and are unique by
  \ref{M:ind}(b--d)), then
  $\semr{\case{I,\tup{\tau},Q}{t}{\tup{u}}}{\rho} =
   \semr{u_j}{\rho}(\tup{d})$ (iterated application); if the case
  term is a proof its denotation is $\emptyset$ by \ref{M:typing}.
\item \label{M:params} \emph{(Parameters)}
  For every parameter declaration $c : T$ in $\Env$:
  $\sem{c} \in \sem{T}$.
\end{enumerate}
\end{definition}

\begin{definition}[Admissible environment; validity]
\label{def:admissible}
An environment $\Env$ is \emph{admissible} if a standard model of it
exists.  Fix an admissible $\Env$ and a standard model $\MM$.  A
closed proposition $\varphi$ is \emph{valid} (in $\MM$) if
$\sem{\varphi} = \holds$.
\end{definition}

Note that admissibility already builds in the truth of all premises:
if $c : \varphi$ is a parameter declaration with $\varphi$ a closed
proposition, then \ref{M:params} gives $\sem{c} \in \sem{\varphi} \in
\Om$, hence $\sem{\varphi} = \holds$, i.e.\ $\varphi$ is valid.  The
same holds for anything \emph{derivable}:

\begin{proposition}[Provability implies validity]
\label{prop:prov-valid}
If $\vdash t : \varphi$ for a closed proposition $\varphi$, then
$\varphi$ is valid.
\end{proposition}

\begin{proof}
By \ref{M:typing}, $\sem{t} \in \sem{\varphi}$, so $\sem{\varphi}
\neq \emptyset$; since $\sem{\varphi} \in \sem{\SProp} = \Om$ by
\ref{M:typing} applied to $\varphi : \SProp$, we get $\sem{\varphi} =
\holds$.
\end{proof}

\begin{proposition}[Existence]
\label{prop:model-exists}
Every environment arising from a consistent development in the
standard set-theoretic interpretation of CIC is admissible.  More
precisely, the proof-irrelevant set-theoretic model constructions of
Werner~\cite{Werner1997}, Aczel~\cite{Aczel1998} and
Lee--Werner~\cite{LeeWerner2011} interpret a calculus containing
$\CICm$ in ZFC with inaccessibles and satisfy
\ref{M:typing}--\ref{M:params}; in particular, environments
consisting of well-typed definitions, data inductive declarations,
and parameter declarations whose types are interpretable (e.g.\
premises that are provable, or axioms true in the intended
interpretation, such as functional extensionality, proof
irrelevance, excluded middle or classical choice) are admissible.
\end{proposition}

\begin{proof}[Proof notes]
We do not reproduce the model construction; we indicate why each
clause of \cref{def:model} holds in it.  Products into $\SProp$ are
interpreted by the trace encoding exactly as in \ref{M:pi} (this is
Werner's proof-irrelevant interpretation; \cite{Werner1997,
LeeWerner2011}); products into type universes are set-theoretic
dependent products, with $\Uone$ closed under them by inaccessibility
of $\kappa_1$ (for a product over a proposition domain into a type
universe there is nothing to check --- such products are not
well-formed in $\CICm$).  \ref{M:conv} and \ref{M:subst} are the
standard soundness lemmas of the interpretation.  Connectives and the
primitive equality \ref{M:conn} match the interpretation of the
corresponding inductive definitions of CIC under proof irrelevance:
e.g.\ Coq's \texttt{eq} is interpreted as a family that is $\holds$
on the diagonal and $\emptyset$ off it.  Inductive data types are
interpreted as least fixed points of their constructor rules with
constructors modelled by injectively tagged tuples, giving
\ref{M:ind}(a--d); $\iota$-reduction soundness gives \ref{M:case}.
\ref{M:params} is the definition of interpretability of the
parameters.
\end{proof}

\begin{remark}
\label{rem:model-choice}
Everything that follows is relative to one fixed standard model
$\MM$ of the ambient environment.  Which standard model is chosen is
irrelevant to the theorems; what matters is that all premises given
to the hammer are valid in it, which admissibility guarantees.  Note
also a convenient consequence of \ref{M:ind}: concrete data types
have their expected denotations.  For instance the empty inductive
type (no constructors) denotes $\emptyset$; a one-constructor,
zero-argument type ($\mathsf{unit}$) denotes a singleton;
$\sem{\mathsf{nat}}$ is Dedekind-infinite, since
$\sem{\mathsf{S}}$ is injective by \ref{M:ind}(b) and omits
$\sem{\mathsf{O}}$ from its image by \ref{M:ind}(c); and
$\sem{\mathsf{list}}(a) $ is infinite whenever $a \neq \emptyset$,
finite (a singleton) when $a = \emptyset$.  These facts drive the
counterexamples and the infinity criterion of \cref{sec:erasure}.
\end{remark}
